% =====================================================================
%  Цель и задание лабораторной работы.
% =====================================================================

\textbf{Цель:} оценить защищённость самостоятельно разработанного
приложения -- проанализировать политику безопасности операционной
системы, компонентов приложения и их взаимодействия, сетевую политику
безопасности и защищённость конфиденциальной информации -- и получить
практические навыки обоснованной защиты пользовательских данных на
уровне архитектуры приложения, а не только на уровне ОС или СУБД.

\vspace{1em}

\textbf{Задание}:
\begin{enumerate}
  \item Разработать приложение, которое позволяет:
  \begin{itemize}
    \item добавить пользователя в систему;
    \item авторизовать пользователя на основе идентификационных данных;
    \item создать, редактировать, удалить, осуществить поиск
    конфиденциальных данных;
    \item создать, редактировать, удалить, осуществить поиск
    неконфиденциальных данных;
    \item деавторизовать авторизованного пользователя.
  \end{itemize}
  \item Выполнить анализ разработанного приложения с точки зрения:
  политики безопасности операционной системы; политики безопасности
  компонентов приложения и их взаимодействия между собой; сетевой
  политики безопасности; безопасности конфиденциальности данных.
  \item Оформить отчёт о проделанной работе.
\end{enumerate}

\textbf{Постановка задачи.} Приложение \texttt{SecureApp} реализовано как
типовое клиент-серверное веб-приложение: SPA-клиент на React обращается по
HTTP(S) к REST API на ASP.NET Core, который хранит данные в локальной базе
SQLite. Приложение ведёт учёт пользователей и двух типов пользовательских
записей -- конфиденциальных и неконфиденциальных -- в объёме, достаточном
для содержательного анализа всех четырёх аспектов защищённости,
перечисленных в задании, без привязки к какой-либо предметной области.
